\section{Igualdad organizacional: por qué la investigación sobre Agent necesita diversidad}

En la entrevista, la «igualdad organizacional» no es una consigna política, sino un método para organizar la investigación. Luo Fuli (罗福莉) considera que el posentrenamiento necesita diversity, y que la participación de quienes trabajan en preentrenamiento es un buen complemento. Un sistema de Agent abarca el modelo, los datos, el RL, el producto, la ingeniería, el costo, la seguridad y el entorno del usuario; para un equipo de un solo perfil es difícil cubrir todos estos problemas.

\begin{figure}[H]
\centering
\includegraphics[width=0.92\textwidth]{org-equality.png}
\caption{Igualdad organizacional en la investigación sobre Agent: varias funciones entran en el mismo ciclo cerrado de investigación.}
\end{figure}

\begin{knowledgebox}{Lectura de la figura: la igualdad organizacional no es igualitarismo}
En la figura, pre-train, post-train/RL, product/UX, infra/systems y leadership se conectan con Agent research. Esto indica que las distintas funciones deben entrar en el ciclo cerrado de investigación, en lugar de que un único grupo monopolice el juicio. La igualdad significa que la información y la responsabilidad entran en un mismo sistema, no que todos hagan el mismo trabajo.
\end{knowledgebox}

\subsection{Cómo la inteligencia colectiva mejora el framework de Agent}

Durante el Año Nuevo chino, Luo Fuli impulsó que todo el equipo usara OpenClaw. Ella subraya que la imaginación de una sola persona es limitada; cuando los miembros del equipo muestran en el grupo general «resulta que también se puede usar así», despiertan la imaginación de los demás. La inteligencia colectiva no es un concepto abstracto: acelera la investigación al compartir experiencias de uso, estimular mutuamente la imaginación de escenarios y descubrir con rapidez las capacidades y las carencias del framework.

\begin{importantbox}{La organización también es un entorno de entrenamiento}
Si un Agent puede aprender de su entorno, un equipo también puede aprender de un entorno de uso compartido. Hacer que el equipo entre en conjunto a nuevas herramientas, nuevos paradigmas y nuevos entornos de tareas equivale, en esencia, a aplicarle un post-train a la organización.
\end{importantbox}

\subsection{Cultura y valores}

La descripción del video subraya que, en este cambio tecnológico profundo, el fundamento de un AI Lab son su cultura y sus valores. Más adelante, Luo Fuli también comenta que lo que hace cada día debe mejorar un poco el mundo: que el trabajo tedioso se sustituya y que las personas tengan tiempo para hacer cosas de mayor valor. Estos valores influyen en cómo el equipo elige sus tareas, cómo define sus benchmarks y cómo asigna sus recursos.

\begin{warningbox}{Los valores no son un adorno}
A medida que los modelos pueden sustituir cada vez más el trabajo humano, lo que un equipo decide optimizar, sustituir o conservar se vuelve parte de las decisiones técnicas. Sin valores que lo acoten, cuanto más capaz sea un Agent, mayor puede ser también el riesgo.
\end{warningbox}

\subsection{Resumen de la sección}

En la era de los Agent, la capacidad organizacional no consiste solo en contratar a más personas, sino en hacer que personas con distintos perfiles entren juntas en el ciclo cerrado de investigación. La igualdad organizacional, la inteligencia colectiva y los valores determinan si un laboratorio puede adoptar con rapidez un nuevo paradigma.
